% Part: computability
% Chapter: computability-theory
% Section: applications-fixed-point

\documentclass[../../../include/open-logic-section]{subfiles}

\begin{document}

\olfileid[tr]{cmp}{thy}{apf}
\olsection{Sabit Nokta Teoreminin Uygulanması}

Sabit nokta teoremi özünde kısmi hesaplanabilir fonksiyonları indisleri
cinsinden tanımlamamızı sağlar. Örneğin her $y$ için
\[
  \cfind{e}(y)=e+y
\]
olacak bir $e$ indisi bulabiliriz. Başka bir örnek olarak sabit nokta
teoreminin ispatı, Java ya da C++ ile kendisini çıktı olarak veren bir program
tasarlamak için kullanılabilir.

Her $e$ için $W_e$, $\cfind{e}$ fonksiyonunun tanım kümesi olursa
$W_0,W_1,W_2,\dots$ dizisinin hesaplanabilir biçimde numaralandırılabilir
kümeleri numaralandırdığını anımsayın. Bu kümelerin bazıları
hesaplanabilirdir. Girdi olarak bir $x$ değeri alan ve $W_x$ hesaplanabilir
olduğunda onun karakteristik fonksiyonunun bir indisini döndüren bir
algoritmanın bulunup bulunmadığı sorulabilir. Yanıt ``hayır''dır; böyle bir
algoritma yoktur.

\begin{thm}
Şu özelliğe sahip hiçbir kısmi hesaplanabilir $f$ fonksiyonu yoktur:
$W_e$ hesaplanabilir olduğunda $f(e)$ tanımlı ve $\cfind{f(e)}$, $W_e$'nin
karakteristik fonksiyonu olsun.
\end{thm}

\begin{proof}
$f$ herhangi bir hesaplanabilir fonksiyon olsun. $W_e$ hesaplanabilir,
fakat $\cfind{f(e)}$ onun karakteristik fonksiyonu olmayacak bir $e$
kuracağız. Sabit nokta teoremini kullanarak
\[
  \cfind{e}(y)\simeq\begin{cases}
    0 & \text{$y=0$ ve $\cfind{f(e)}(0)\fdefined=0$ ise}\\
    \text{tanımsız} & \text{aksi durumda}
  \end{cases}
\]
olacak bir $e$ indisi bulabiliriz. Yani $e$, sabit nokta teoreminin
\[
  g(x,y)\simeq\begin{cases}
    0 & \text{$y=0$ ve $\cfind{f(x)}(0)\fdefined=0$ ise}\\
    \text{tanımsız} & \text{aksi durumda}
  \end{cases}
\]
fonksiyonuna uygulanmasıyla elde edilir.
Biçimsel olmayan biçimde $g$'nin kısmi hesaplanabilir olduğunu şöyle
görebiliriz: Algoritma $x,y$ girdisinde önce $y=0$ olup olmadığını denetler.
Öyleyse $f(x)$'i, ardından evrensel makineyle $\cfind{f(x)}(0)$ değerini
hesaplar. Bu son hesap durup $0$ döndürürse algoritma $0$ döndürür; aksi
durumda durmaz.

Şimdi $\cfind{f(e)}(0)$ tanımlı ve $0$'a eşitse $\cfind{e}(y)$ tam olarak
$y=0$ olduğunda tanımlıdır; dolayısıyla $W_e=\{0\}$ olur.
$\cfind{f(e)}(0)$ tanımsızsa ya da tanımlı fakat $0$'a eşit değilse
$W_e=\emptyset$ olur. Her iki durumda da $\cfind{f(e)}$, $W_e$'nin
karakteristik fonksiyonu değildir; çünkü $0$ girdisinde yanlış yanıt verir.
\end{proof}

\end{document}

